% =====================================================================
%  CHAPTER 1: 绪论（对应大纲第一章：绪论）
%  参考PPT：第一章 绪言（郝加虎，安徽医科大学）
% =====================================================================
\chapter{绪论}
\setcounter{chapter}{1}
\chapterintro{
\textbf{掌握：}儿童少年生长发育的规律及其卫生学意义。\newline
\textbf{熟悉：}课程的性质、任务、特点、重要性、研究内容和研究方法。
}

% ----- 1.1 儿童少年卫生学概述 -----
\section{儿童少年卫生学概述}

\subsection{研究对象}

\begin{defbox}
\textbf{儿童少年卫生学（child and adolescent health）：}研究维护和促进儿童少年健康的一门学科。研究儿童少年身心发育随年龄变化的特征，分析影响生长发育的遗传和环境因素，阐明儿童少年机体与学习及生活环境间的相互关系，制定相应卫生要求和措施，预防疾病、增强体质，促进身心健康发育，并为成年期健康奠定良好基础，从而达到提高生命质量的目的。

\textbf{研究对象：}0-24岁儿童少年；重点人群：中小学生群体。
\end{defbox}

\subsection{儿童少年年龄范围的界定}

\begin{keybox}
\textbf{儿童少年年龄分期：}

\begin{table}[H]
\centering
\caption{儿童少年年龄分期}
\begin{tabularx}{\textwidth}{l>{\raggedright\arraybackslash}p{3.5cm}X}
\hline
\textbf{分期} & \textbf{英文名称} & \textbf{年龄范围} \\
\hline
\imp{胎儿期} & fetal period & 从受精卵形成开始到孕40周胎儿娩出 \\
\imp{婴儿期} & infant period & 出生后1年 \\
\imp{新生儿期} & neonatal period & 出生后28天（属于婴儿期） \\
\imp{幼儿期} & toddle period & 出生后第2-3年 \\
\imp{学龄前期} & preschool period & 3-6岁 \\
\imp{学龄期} & school-age period & 6-11/12岁，相当于小学阶段 \\
\imp{青春期} & adolescence & 从青春发动开始到生长基本成熟（10-19岁） \\
\imp{青年期} & youth & 15-24岁 \\
\hline
\end{tabularx}
\end{table}
\end{keybox}

\subsection{儿童少年生物和社会特征}

\begin{keybox}
\textbf{儿童少年三大生物和社会特征：}

\begin{enumerate}[leftmargin=*]
    \item \imp{处于生长发育过程中}——生长发育是儿童少年的基本特征之一，也是社会发展的一面镜子。了解儿童少年生长发育规律，并深入研究其影响因素，是制定儿童少年保健工作的相应政策和行动纲领的重要前提。

    \item \imp{接受学校教育的塑造}——接受教育是儿童少年的一个重要社会特征，学校集体生活构成其重要的生活内容。创造良好的学校物质环境和社会支持环境，建立健康促进学校，开展综合性学校卫生服务，是促进其身心健康和生长发育潜能的重要条件。

    \item \imp{具有社会脆弱性和健康易损性}——在全人口中，儿童少年由于其身心以及各种能力发育的不完善，始终是一个弱势的群体。儿童神经系统发育不成熟，与成人相比，更容易受到有害生物和环境因素影响，从而造成发育和行为异常。在不同的社会生态环境层面——个体、人际间、社区、结构和政策层面，以及在不同的发展阶段，儿童少年的心理健康受到多个动态因素的影响。
\end{enumerate}
\end{keybox}

\subsection{儿童少年生存与发展权利}

\begin{defbox}
\textbf{儿童少年生存与发展权利：}

\begin{enumerate}[leftmargin=*]
    \item \imp{生存权}：每一个儿童都应该享有的对自己生命权以及维持基本健康生活需求的生活条件保障权。

    \item \imp{发展权}：保障儿童成长过程中的各种需要得到满足，包括儿童有接受一切形式的教育（正规的和非正规的教育）的权利，能够给予儿童的身体、心理、精神、道德与社交发展的相应生活水平。发展权包括参与权、受保护权。
\end{enumerate}
\end{defbox}

% ----- 1.2 生长发育规律 -----
\section{儿童少年生长发育的规律及其卫生学意义}

\begin{keybox}
\textbf{儿童少年生长发育的规律：}

生长发育是儿童少年的基本特征之一，也是社会发展的一面镜子。儿童少年身心发育随年龄变化的特征，受遗传和环境因素共同影响。生长发育是累积和连续的过程，个体从"胎儿—儿童—少年—青壮年—老年"构成完整的生命历程，每一个阶段健康既是前一阶段的延续，也是下一阶段的基础。

\textbf{卫生学意义：}

\begin{enumerate}[leftmargin=*]
    \item 了解儿童少年生长发育规律，并深入研究其影响因素，是制定儿童少年保健工作的相应政策和行动纲领的重要前提。
    \item 儿童少年发展潜力的发挥必须通过整个生命历程中的支持性环境来实现。"支持性环境"是能够激发儿童少年最佳发展潜力的系统、力量和因素，是由贯穿整个生命历程的卫生服务、政策、计划和社区共同参与。
    \item 生长发育规律为制定相应卫生要求和措施、预防疾病、增强体质、促进身心健康发育提供科学依据，并为成年期健康奠定良好基础，从而达到提高生命质量的目的。
\end{enumerate}
\end{keybox}

% ----- 1.3 儿童少年卫生学学科体系 -----
\section{儿童少年卫生学学科体系}

\subsection{研究内容}

\begin{defbox}
\textbf{儿童少年卫生学的四大研究内容：}

\begin{enumerate}[leftmargin=*]
    \item \imp{儿童少年生长发育规律}——生长发育作为人类最为复杂的生物学现象，是本学科研究的永恒主题，也是学科首先要研究和认识的理论问题。包括身体发育和心理发育两个方面。

    \item \imp{儿童少年健康状况及其决定因素}——以保护、促进、增强儿童少年身心健康为宗旨，提出相应的卫生要求和适宜的卫生措施，维护儿童少年的健康，减少疾病，预防死亡，提高生存质量。

    \item \imp{儿童少年卫生服务}——包括：生长发育和健康监测；健康教育和健康促进；常见病防治与健康管理；心理卫生服务；学校营养服务；校园安全管理与伤害及暴力预防；体育与体力活动促进。

    \item \imp{儿童少年卫生学的技术和方法}——人体测量、健康相关体能测试、行为评定等。
\end{enumerate}
\end{defbox}

\subsection{三级预防理论}

\begin{keybox}
\textbf{三级预防理论：}

三级预防理论是学科发展的重要理论构架。三级预防是儿童少年健康促进的首要和有效手段，是现代医学为人们提供的健康保障。

\begin{enumerate}[leftmargin=*]
    \item \imp{一级预防（primary prevention）/ 病因预防}：疾病尚未发生时针对致病因素或危险因素采取措施，是预防和消灭疾病的根本措施。是最积极最有效的预防措施。

    \item \imp{二级预防（secondary prevention）/ "三早"预防}：早发现、早诊断、早治疗，是防止或减缓疾病发展采取的措施。

    \item \imp{三级预防（tertiary prevention）/ 临床预防}：防治伤残和促进功能恢复，提高生命质量，延长寿命，降低病死率，主要是对症治疗和康复治疗措施。
\end{enumerate}
\end{keybox}

% ----- 复习题 -----
\section{复习题}

\begin{enumerate}[leftmargin=*, itemsep=8pt]
    \item \textbf{【名词解释】}儿童少年卫生学；新生儿期；青春期；青年期
    \item \textbf{【简答题】}简述儿童少年的年龄分期及各期年龄范围。
    \item \textbf{【简答题】}简述儿童少年的三大生物和社会特征。
    \item \textbf{【简答题】}简述儿童少年的生存与发展权利。
    \item \textbf{【简答题】}简述儿童少年卫生学的四大研究内容。
    \item \textbf{【简答题】}简述三级预防策略的主要内容。
    \item \textbf{【简答题】}试述儿童少年生长发育的规律及其卫生学意义。
\end{enumerate}

\subsection*{参考答案要点}

\begin{enumerate}[leftmargin=*, itemsep=5pt]
    \item \textbf{儿童少年卫生学}：研究维护和促进儿童少年健康的一门学科，研究儿童少年身心发育随年龄变化的特征，分析影响生长发育的遗传和环境因素，阐明儿童少年机体与学习及生活环境间的相互关系，制定相应卫生要求和措施，预防疾病、增强体质，促进身心健康发育，并为成年期健康奠定良好基础，从而达到提高生命质量的目的。研究对象为0-24岁儿童少年，重点人群为中小学生群体。\textbf{新生儿期}：出生后28天，属于婴儿期。\textbf{青春期}：从青春发动开始到生长基本成熟，10-19岁。\textbf{青年期}：15-24岁。

    \item 八个分期：胎儿期（受精卵形成至孕40周娩出）、婴儿期（出生后1年）、新生儿期（出生后28天，属婴儿期）、幼儿期（出生后第2-3年）、学龄前期（3-6岁）、学龄期（6-11/12岁，相当于小学阶段）、青春期（10-19岁，从青春发动开始到生长基本成熟）、青年期（15-24岁）。

    \item (1) 处于生长发育过程中——生长发育是儿童少年的基本特征之一，也是社会发展的一面镜子；(2) 接受学校教育的塑造——学校集体生活构成其重要的生活内容；(3) 具有社会脆弱性和健康易损性——身心及各种能力发育不完善，始终是弱势群体，更容易受到有害生物和环境因素影响。

    \item (1) 生存权：每一个儿童都应该享有的对自己生命权以及维持基本健康生活需求的生活条件保障权；(2) 发展权：保障儿童成长过程中的各种需要得到满足，包括接受一切形式教育的权利，给予儿童身体、心理、精神、道德与社交发展的相应生活水平。发展权包括参与权、受保护权。

    \item (1) 儿童少年生长发育规律——学科永恒主题，包括身体发育和心理发育；(2) 儿童少年健康状况及其决定因素——以促进身心健康为宗旨，提出卫生要求和措施；(3) 儿童少年卫生服务——生长发育和健康监测、健康教育和健康促进、常见病防治与健康管理、心理卫生服务、学校营养服务、校园安全管理与伤害及暴力预防、体育与体力活动促进；(4) 儿童少年卫生学的技术和方法——人体测量、健康相关体能测试、行为评定等。

    \item 一级预防（病因预防）：疾病尚未发生时针对致病因素或危险因素采取措施，是预防和消灭疾病的根本措施，是最积极最有效的预防措施；二级预防（"三早"预防）：早发现、早诊断、早治疗，是防止或减缓疾病发展采取的措施；三级预防（临床预防）：防治伤残和促进功能恢复，提高生命质量，延长寿命，降低病死率，主要是对症治疗和康复治疗措施。

    \item 生长发育是儿童少年的基本特征，身心发育随年龄变化，受遗传和环境因素共同影响。生长发育是累积和连续的过程，构成从胎儿到老年的完整生命历程，每一阶段健康是前一阶段的延续和下一阶段的基础。其卫生学意义：(1) 了解生长发育规律并研究其影响因素，是制定儿童少年保健工作政策和行动纲领的重要前提；(2) 儿童少年发展潜力的发挥需通过整个生命历程中的支持性环境实现；(3) 生长发育规律为制定卫生要求和措施、预防疾病、增强体质、促进身心健康发育提供科学依据，为成年期健康奠定基础，提高生命质量。
\end{enumerate}

\newpage
